\documentclass[11pt,a4paper]{article}
\usepackage[utf8]{inputenc}
\usepackage[french]{babel}
\usepackage[T1]{fontenc}
\usepackage{geometry}
\usepackage{array}
\usepackage{amsmath}
\usepackage{enumitem}
\usepackage{tikz-osci}
\usetikzlibrary{babel}

\geometry{a4paper, margin=2cm}

\begin{document}

% En-tête du contrôle
\noindent
\begin{tabular}{|p{5cm}|p{5cm}|p{5cm}|}
\hline
\textbf{Mercredi 20 octobre} & \textbf{Sciences -- Contrôle 1} & \textbf{Classe de TSN1} \\
\hline
\textbf{Nom :} & & \textbf{Note :} \\
 & & \\
\textbf{Prénom :} & & \\
\hline
\end{tabular}

\vspace{1em}

\section*{Exercice 1 :}
Les oscillogrammes ci-dessous ont été obtenus avec le même balayage (5 ms/div) et la même sensibilité (2 V/div).

\vspace{0.5em}

\begin{center}
\begin{minipage}{0.18\textwidth}\centering\textbf{oscillogramme n°1}\\[0.2cm]
\osci[scale=0.25,second channel=0,screen offset one=0,time div=5,voltage div one=2,sample rate=200,xy mode=0,func one=4*sin(2*pi*50*x),color one=009900,graph back color=FFFFFF,info back color=D6D6D6,info text color=000000,main axis color=000000,grid color=CCCCCC]
\end{minipage}
\begin{minipage}{0.18\textwidth}\centering\textbf{oscillogramme n°2}\\[0.2cm]
\osci[scale=0.25,second channel=0,screen offset one=0,time div=5,voltage div one=2,sample rate=200,xy mode=0,func one=5,color one=009900,graph back color=FFFFFF,info back color=D6D6D6,info text color=000000,main axis color=000000,grid color=CCCCCC]
\end{minipage}
\begin{minipage}{0.18\textwidth}\centering\textbf{oscillogramme n°3}\\[0.2cm]
\osci[scale=0.25,second channel=0,screen offset one=0,time div=5,voltage div one=2,sample rate=200,xy mode=0,func one=6*sin(2*pi*50*x),color one=009900,graph back color=FFFFFF,info back color=D6D6D6,info text color=000000,main axis color=000000,grid color=CCCCCC]
\end{minipage}
\begin{minipage}{0.18\textwidth}\centering\textbf{oscillogramme n°4}\\[0.2cm]
\osci[scale=0.25,second channel=0,screen offset one=0,time div=5,voltage div one=2,sample rate=200,xy mode=0,func one=6*abs(sin(2*pi*50*x)),color one=009900,graph back color=FFFFFF,info back color=D6D6D6,info text color=000000,main axis color=000000,grid color=CCCCCC]
\end{minipage}
\begin{minipage}{0.18\textwidth}\centering\textbf{oscillogramme n°5}\\[0.2cm]
\osci[scale=0.25,second channel=0,screen offset one=0,time div=5,voltage div one=2,sample rate=200,xy mode=0,func one=(sin(2*pi*250*x)>0)*6,color one=009900,graph back color=FFFFFF,info back color=D6D6D6,info text color=000000,main axis color=000000,grid color=CCCCCC]
\end{minipage}
\end{center}

\vspace{1em}

\begin{enumerate}
\item Quelles sont les tensions qui ont la même période ? Calculer alors cette période et la fréquence correspondante.

\vspace{3cm}

\item Quels sont les oscillogrammes qui ont la même tension maximale ? Calculer alors sa valeur.

\vspace{3cm}

\item Pour l'oscillogramme numéro 3, calculer la tension maximale, puis la tension efficace.

\vspace{3cm}

\end{enumerate}

\begin{minipage}[t]{0.48\textwidth}
\begin{enumerate}
\setcounter{enumi}{3}
\item Si pour l'oscillogramme numéro 3 on avait réglé la sensibilité à 1 V/div, quel aurait été son aspect ?

\vspace{0.3cm}

\begin{center}
% Grille vide pour la réponse de l'élève
\osci[%
scale=0.75,
second channel=0,
screen offset one=0,
time div=5,
voltage div one=1,
sample rate=200,
xy mode=0,
func one=0,
color one=009900,
graph back color=FFFFFF,
info back color=D6D6D6,
info text color=000000,
main axis color=000000,
grid color=CCCCCC
]
\end{center}
\end{enumerate}
\end{minipage}
\hfill
\begin{minipage}[t]{0.48\textwidth}
\begin{enumerate}
\setcounter{enumi}{4}
\item Si pour l'oscillogramme n°3 on avait réglé le balayage à 10 ms/div (en remettant la sensibilité à 2 V/div), quel aurait été son aspect ?

\vspace{0.3cm}

\begin{center}
% Grille vide pour la réponse de l'élève
\osci[%
scale=0.75,
second channel=0,
screen offset one=0,
time div=10,
voltage div one=2,
sample rate=200,
xy mode=0,
func one=0,
color one=009900,
graph back color=FFFFFF,
info back color=D6D6D6,
info text color=000000,
main axis color=000000,
grid color=CCCCCC
]
\end{center}
\end{enumerate}
\end{minipage}

\newpage

\section*{Exercice 2 :} 
QCM (cocher la ou les bonnes réponses)

\vspace{0.5em}

\small
\begin{enumerate}
\item Le redressement à l'aide d'un pont de diodes permet:
\begin{itemize}[label=$\square$]
\item de fournir une tension continue à partir d'une tension alternative
\item d'obtenir un interrupteur dans le circuit
\item de fournir une tension alternative à partir d'une tension continue
\end{itemize}

\item En observant la diode, le courant circule de :
\begin{itemize}[label=$\square$]
\item la droite vers la gauche
\item la gauche vers la droite
\item dans les deux sens
\end{itemize}

\item Une diode est un composant:
\begin{itemize}[label=$\square$]
\item polarisé
\item linéaire
\item possédant trois bornes
\end{itemize}

\item Le chargeur délivre:
\begin{itemize}[label=$\square$]
\item du courant continu
\item du courant alternatif
\item de l'énergie chimique
\end{itemize}

\item Il est possible de recharger:
\begin{itemize}[label=$\square$]
\item uniquement des accumulateurs
\item les accumulateurs et les piles
\item uniquement les piles
\end{itemize}

\item Un accumulateur délivre un courant :
\begin{itemize}[label=$\square$]
\item alternatif
\item continu
\item redressé
\end{itemize}

\item Une diode :
\begin{itemize}[label=$\square$]
\item autorise le passage du courant dans les deux sens
\item autorise le passage du courant dans un seul sens
\item n'autorise pas le passage du courant
\end{itemize}

\item Une tension redressée est une tension :
\begin{itemize}[label=$\square$]
\item qui garde la même valeur
\item qui garde le même signe
\item toujours positive
\end{itemize}

\item Lorsque l'on réalise le montage correspondant au schéma ci-dessus et que l'on ferme l'interrupteur :
\begin{itemize}[label=$\square$]
\item Le courant ne circule pas dans le circuit
\item La lampe ne s'allume pas dans le circuit
\item La lampe s'allume dans le circuit
\end{itemize}

\item Une tension redressée par un pont de diode correspond :

\begin{minipage}{0.3\textwidth}
$\square$ 
\begin{center}
\begin{tikzpicture}[scale=0.4]
\draw[domain=0:4*pi,smooth,variable=\x,blue,thick] plot ({\x},{sin(\x r)});
\draw[->] (0,-1.2) -- (0,1.5);
\draw[->] (0,0) -- (13,0);
\end{tikzpicture}
\end{center}
\end{minipage}
\begin{minipage}{0.3\textwidth}
$\square$
\begin{center}
\begin{tikzpicture}[scale=0.4]
\draw[domain=0:4*pi,smooth,variable=\x,blue,thick] plot ({\x},{abs(sin(\x r))});
\draw[->] (0,-0.2) -- (0,1.5);
\draw[->] (0,0) -- (13,0);
\end{tikzpicture}
\end{center}
\end{minipage}
\begin{minipage}{0.3\textwidth}
$\square$
\begin{center}
\begin{tikzpicture}[scale=0.4]
\draw[blue,thick] (0,1) -- (12.5,1);
\draw[->] (0,-0.2) -- (0,1.5);
\draw[->] (0,0) -- (13,0);
\end{tikzpicture}
\end{center}
\end{minipage}

\item La tension obtenue après un chargeur est :

\begin{minipage}{0.3\textwidth}
$\square$ 
\begin{center}
\begin{tikzpicture}[scale=0.4]
\draw[domain=0:4*pi,smooth,variable=\x,blue,thick] plot ({\x},{sin(\x r)});
\draw[->] (0,-1.2) -- (0,1.5);
\draw[->] (0,0) -- (13,0);
\end{tikzpicture}
\end{center}
\end{minipage}
\begin{minipage}{0.3\textwidth}
$\square$
\begin{center}
\begin{tikzpicture}[scale=0.4]
\draw[domain=0:4*pi,smooth,variable=\x,blue,thick] plot ({\x},{abs(sin(\x r))});
\draw[->] (0,-0.2) -- (0,1.5);
\draw[->] (0,0) -- (13,0);
\end{tikzpicture}
\end{center}
\end{minipage}
\begin{minipage}{0.3\textwidth}
$\square$
\begin{center}
\begin{tikzpicture}[scale=0.4]
\draw[blue,thick] (0,1) -- (12.5,1);
\draw[->] (0,-0.2) -- (0,1.5);
\draw[->] (0,0) -- (13,0);
\end{tikzpicture}
\end{center}
\end{minipage}

\end{enumerate}
\normalsize

\vspace{1em}

\section*{Exercice 3 :} 
Pour chacun des schémas suivants, indiquer les lampes qui s'allument.

\vspace{0.5em}

\begin{center}
\textit{[Insérer les 3 schémas de circuits avec lampes et diodes]}
\end{center}

\vspace{0.5cm}

\textbf{Schéma 1 :}

La lampe L$_1$ est \dotfill

La lampe L$_2$ est \dotfill

\vspace{1cm}

\textbf{Schéma 2 :}

La lampe L$_1$ est \dotfill

La lampe L$_2$ est \dotfill

La lampe L$_3$ est \dotfill

\vspace{1cm}

\textbf{Schéma 3 :}

La lampe L$_1$ est \dotfill

La lampe L$_2$ est \dotfill

La lampe L$_3$ est \dotfill

\newpage

\section*{Exercice 4 :}
Le schéma suivant représente une partie d'une alimentation d'antenne parabolique de télévision.

\vspace{0.5cm}

\begin{center}
\textit{[Insérer le schéma du circuit avec pont de diodes, condensateur C, résistance R et générateur G]}
\end{center}

\begin{center}
\begin{minipage}{0.45\textwidth}\centering\textbf{Oscillogramme 1}\\[0.2cm]
\osci[scale=0.7,second channel=0,screen offset one=0,time div=5,voltage div one=2,sample rate=200,xy mode=0,func one=15*sin(2*pi*50*x),color one=009900,graph back color=FFFFFF,info back color=D6D6D6,info text color=000000,main axis color=000000,grid color=CCCCCC]
\end{minipage}
\hfill
\begin{minipage}{0.45\textwidth}\centering\textbf{Oscillogramme 2}\\[0.2cm]
\osci[scale=0.7,second channel=0,screen offset one=0,time div=5,voltage div one=2,sample rate=200,xy mode=0,func one=15*abs(sin(2*pi*50*x)),color one=009900,graph back color=FFFFFF,info back color=D6D6D6,info text color=000000,main axis color=000000,grid color=CCCCCC]
\end{minipage}
\end{center}

\vspace{0.5cm}

\textbf{Oscillogramme 3} (après fermeture de K) :

\begin{center}
% Oscillogramme 3 : UR lissée par le condensateur
\osci[%
scale=0.8,
second channel=0,
screen offset one=0,
time div=5,
voltage div one=2,
sample rate=200,
xy mode=0,
func one=13.5+0.5*sin(2*pi*100*x),
color one=009900,
graph back color=FFFFFF,
info back color=D6D6D6,
info text color=000000,
main axis color=000000,
grid color=CCCCCC
]
\end{center}

\vspace{1em}

\begin{enumerate}
\item L'interrupteur K étant ouvert. À l'aide d'un oscilloscope, on obtient les oscillogrammes représentés ci-dessus :
\begin{itemize}[label=--]
\item oscillogramme 1 : tension U$_G$ délivrée par le générateur
\item oscillogramme 2 : tension U$_R$ aux bornes de la résistance R.
\end{itemize}

\begin{enumerate}
\item Quel est le rôle du pont de diodes, appelé "Pont de Graëtz" ?

\dotfill

\dotfill

\item À l'aide de l'oscillogramme 1, \textbf{déterminer} la période, la fréquence puis la valeur maximale U$_{Gmax}$ de la tension délivrée par le générateur.

\dotfill

\dotfill

\dotfill
\end{enumerate}

\item A l'aide de l'oscillogramme 2,

\begin{enumerate}
\item \textbf{Déterminer} la période, la fréquence puis la valeur maximale U$_{Rmax}$ de la tension visualisée aux bornes de la résistance.

\dotfill

\dotfill

\dotfill

\item \textbf{Entourer} quelles sont les deux valeurs qui ne peuvent pas être lues sur le voltmètre (en dérivation par rapport à la résistance) parmi les trois valeurs suivantes : 0 V ; 15 V ; 10,6 V.
\end{enumerate}

\item On ferme l'interrupteur K.

À l'aide de l'oscilloscope, on obtient l'oscillogramme 3, représentant la nouvelle tension U$_R$ aux bornes de la résistance.

\begin{enumerate}
\item Quel est le rôle du condensateur ?

\dotfill

\dotfill

\item Le voltmètre indique une tension $U_R = 13{,}5$ V lorsque la résistance est $R = 10\,\Omega$. \textbf{Calculer} l'intensité du courant traversant la résistance R.

\vspace{3cm}

\end{enumerate}

\item Flécher le sens du courant sur le schéma du montage lors des deux alternances, un schéma pour chaque alternance.

\vspace{0.5cm}

\begin{center}
\textit{[Insérer les deux schémas du pont de diodes avec flèches indiquant le sens du courant]}
\end{center}

\vspace{3cm}

\end{enumerate}

\end{document}